\chapter{Les variables structurées}
\section{Le type structuré}
{
Dans les chapitres précédents, nous avons utilisé des variables de types
dits «~simples~» (entier, réel, booléen, caractère, chaine) ne pouvant
contenir qu’une seule valeur à la fois. Cependant, certains types
d’information consistent en un regroupement de plusieurs données
élémentaires. Quelques exemples :}

\liststyleListv
\begin{itemize}
\item {
une \textbf{date} est composée de trois éléments (le jour, le mois,
l’année), le mois s’exprimant tantôt par un entier (15/10/2004) tantôt
par une chaine (15 octobre 2004)}
\item {
un \textbf{moment} de la journée est un triple d’entiers (heures,
minutes, secondes)}
\item {
la localisation d’un \textbf{point} dans un plan nécessite la
connaissance de deux coordonnées cartésiennes (l’abscisse \textit{x} et
l’ordonnée \textit{y}) ou polaires (le rayon \textit{r} et l’angle
\textit{$\theta $})}
\item {
une \textbf{adresse }est composée de plusieurs données : un nom de rue,
un numéro de maison, parfois un numéro de boite postale, un code
postal, le nom de la localité, un nom ou code de pays pour une adresse
à l’étranger…}
\end{itemize}
{
Pour stocker et manipuler de tels ensembles de données, nous utiliserons
des \textbf{types structurés} ou \textbf{structures}\footnote{Nous
verrons dans un chapitre ultérieur que l'approche
«~orienté objet~» offre également une solution élégante pour manipuler
des données complexes}. Une \textbf{structure} est donc un ensemble
fini d’éléments pouvant être de types distincts. Chaque élément de cet
ensemble, appelé \textbf{champ} de la structure, possède un nom unique.
}

{
Noter qu’un champ d’une structure peut lui-même être une structure. Par
exemple, une \textbf{carte d’identité} inclut parmi ses informations
une \textbf{date} de naissance, l’\textbf{adresse} de son
propriétaire…}


\bigskip


\bigskip


\bigskip

\section[Définition d'une structure]{Définition
d'une structure}
{
La définition d’un type structuré adoptera le modèle suivant :}

\begin{center}
 [Warning: Image ignored] % Unhandled or unsupported graphics:
%\includegraphics[width=1.129cm,height=1.282cm]{log1-img/log1-img80}

\end{center}
{\sffamily
\textstyleMotCl{\textcolor{black}{structure}}\textcolor{black}{
\ NomDeLaStructure}}

{\sffamily
\ \ nomChamp1 : type1}

{\sffamily
\ \ nomChamp2 : type2}

{\sffamily
…}

{\sffamily
\ \ nomChampN : typeN}

{\sffamily
\textstyleMotCl{fin} \textstyleMotCl{structure}}

{
\textstylePolicepardfaut{\textcolor{black}{nomChamp1, …, nomChampN sont
les noms des différents champs de la structure, et type1, … , typeN les
types de ces champs. Ces types sont soit les types «~simples~» étudiés
précédemment (entier, réel, booléen, caractère, chaine) soit d’autres
types structurés dont la structure aura été préalablement définie.}}}

{
\textstylePolicepardfaut{Pour exemple, nous définissons ci-dessous trois
types structurés que nous utiliserons souvent par la suite :}}

{\sffamily
\textstyleMotCl{\textcolor{black}{structure}}\textstylePolicepardfaut{\textsf{\textcolor{black}{
Date}}}}

{\sffamily
\textstylePolicepardfaut{\textsf{\ \ jour : entier}}}

{\sffamily
\textstylePolicepardfaut{\textsf{\ \ mois : entier}}}

{\sffamily
\textstylePolicepardfaut{\textsf{\ \ année : entier}}}

{\sffamily
\textstyleMotCl{fin}\textstylePolicepardfaut{\textsf{
}}\textstyleMotCl{structure}}


\bigskip

{\sffamily
\textstyleMotCl{\textcolor{black}{structure}}\textstylePolicepardfaut{\textsf{\textcolor{black}{
Moment}}}}

{\sffamily
\textstylePolicepardfaut{\textsf{\ \ heure : entier}}}

{\sffamily
\textstylePolicepardfaut{\textsf{\ \ minute : entier}}}

{\sffamily
\textstylePolicepardfaut{\textsf{\ \ seconde : entier}}}

{\sffamily
\textstyleMotCl{fin}\textstylePolicepardfaut{\textsf{
}}\textstyleMotCl{structure}}


\bigskip

{\sffamily
\textstyleMotCl{\textcolor{black}{structure}}\textstylePolicepardfaut{\textsf{\textcolor{black}{
Point}}}}

{\sffamily
\textstylePolicepardfaut{\textsf{\ \ x : réel}}}

{\sffamily
\textstylePolicepardfaut{\textsf{\ \ y : réel}}}

{\sffamily
\textstyleMotCl{fin}\textstylePolicepardfaut{\textsf{
}}\textstyleMotCl{structure}}

\section{Déclaration d’une variable de type structuré}
{
\textstylePolicepardfaut{Une fois un type structuré défini, la
déclaration de variables de ce type est identique à celle des variables
simples. Par exemple :}}

{\sffamily
\textstylePolicepardfaut{\textsf{\textcolor{black}{anniversaire, jourJ :
Date}}}}

{\sffamily
\textstylePolicepardfaut{\textsf{\textcolor{black}{départ, arrivée,
unMoment : Moment}}}}

{\sffamily
\textstylePolicepardfaut{\textsf{\textcolor{black}{a, b, centreGravité :
Point}}}}

\section{Utilisation des variables de type structuré}
{
\textstylePolicepardfaut{En général, pour manipuler une variable
structurée ou en modifier le contenu, il faut agir au niveau de ses
champs en utilisant les opérations permises selon leur type. Pour
accéder à l’un des champs d’une variable structurée, il faut mentionner
le nom de ce champ ainsi que celui de la variable dont il fait partie.
Nous utiliserons la notation «~pointée~» :}}


\bigskip

{\sffamily
\textstylePolicepardfaut{\textsf{nomVariable.nomChamp}}}

{
\textstylePolicepardfaut{Exemple d’instructions utilisant les variables
déclarées au paragraphe précédent :}}

{\sffamily
\textstylePolicepardfaut{\textsf{\textcolor{black}{anniversaire}}}\textstylePolicepardfaut{\textsf{\textbf{\textcolor{black}{.}}}}\textstylePolicepardfaut{\textsf{\textcolor{black}{jour
}}}\textstyleMotCl{{\textcolor{black}{←}}}\textstylePolicepardfaut{\textsf{\textcolor{black}{
15}}}}

{\sffamily
\textstylePolicepardfaut{\textsf{\textcolor{black}{anniversaire}}}\textstylePolicepardfaut{\textsf{\textbf{\textcolor{black}{.}}}}\textstylePolicepardfaut{\textsf{\textcolor{black}{mois
}}}\textstyleMotCl{{\textcolor{black}{←}}}\textstylePolicepardfaut{\textsf{\textcolor{black}{
10}}}}

{\sffamily
\textstylePolicepardfaut{\textsf{\textcolor{black}{anniversaire}}}\textstylePolicepardfaut{\textsf{\textbf{\textcolor{black}{.}}}}\textstylePolicepardfaut{\textsf{\textcolor{black}{année
}}}\textstyleMotCl{{\textcolor{black}{←}}}\textstylePolicepardfaut{\textsf{\textcolor{black}{
2004}}}}

{\sffamily
\textstylePolicepardfaut{\textsf{\textcolor{black}{arrivée}}}\textstylePolicepardfaut{\textsf{\textbf{\textcolor{black}{.}}}}\textstylePolicepardfaut{\textsf{\textcolor{black}{heure
}}}\textstyleMotCl{{\textcolor{black}{←}}}\textstylePolicepardfaut{\textsf{\textcolor{black}{
départ}}}\textstylePolicepardfaut{\textsf{\textbf{\textcolor{black}{.}}}}\textstylePolicepardfaut{\textsf{\textcolor{black}{heure
+ 2}}}}

{\sffamily
\textstylePolicepardfaut{\textsf{\textcolor{black}{centreGravité}}}\textstylePolicepardfaut{\textsf{\textbf{\textcolor{black}{.}}}}\textstylePolicepardfaut{\textsf{\textcolor{black}{x
}}}\textstyleMotCl{{\textcolor{black}{←}}}\textstylePolicepardfaut{\textsf{\textcolor{black}{
(a}}}\textstylePolicepardfaut{\textsf{\textbf{\textcolor{black}{.}}}}\textstylePolicepardfaut{\textsf{\textcolor{black}{x
+
b}}}\textstylePolicepardfaut{\textsf{\textbf{\textcolor{black}{.}}}}\textstylePolicepardfaut{\textsf{\textcolor{black}{x)/2}}}}

{
\textstylePolicepardfaut{On peut cependant dans certains cas utiliser
une variable structurée de manière globale (c’est-à-dire d’une façon
qui agit simultanément sur chacun de ses champs). Le cas le plus
courant est l’affectation interne entre deux variables structurées de
même type, par exemple :}}

{\sffamily
\textstylePolicepardfaut{\textsf{\textcolor{black}{arrivée
}}}\textstyleMotCl{{\textcolor{black}{←}}}\textstylePolicepardfaut{\textsf{\textcolor{black}{
départ}}}}

{
\textstylePolicepardfaut{qui résume les trois instructions suivantes :}}

{\sffamily
\textstylePolicepardfaut{\textsf{\textcolor{black}{arrivée}}}\textstylePolicepardfaut{\textsf{\textbf{\textcolor{black}{.}}}}\textstylePolicepardfaut{\textsf{\textcolor{black}{heure
}}}\textstyleMotCl{{\textcolor{black}{←}}}\textstylePolicepardfaut{\textsf{\textcolor{black}{
départ}}}\textstylePolicepardfaut{\textsf{\textbf{\textcolor{black}{.}}}}\textstylePolicepardfaut{\textsf{\textcolor{black}{heure}}}}

{\sffamily
\textstylePolicepardfaut{\textsf{\textcolor{black}{arrivée}}}\textstylePolicepardfaut{\textsf{\textbf{\textcolor{black}{.}}}}\textstylePolicepardfaut{\textsf{\textcolor{black}{minute
}}}\textstyleMotCl{{\textcolor{black}{←}}}\textstylePolicepardfaut{\textsf{\textcolor{black}{
départ}}}\textstylePolicepardfaut{\textsf{\textbf{\textcolor{black}{.}}}}\textstylePolicepardfaut{\textsf{\textcolor{black}{minute}}}}

{\sffamily
\textstylePolicepardfaut{\textsf{\textcolor{black}{arrivée}}}\textstylePolicepardfaut{\textsf{\textbf{\textcolor{black}{.}}}}\textstylePolicepardfaut{\textsf{\textcolor{black}{seconde
}}}\textstyleMotCl{{\textcolor{black}{←}}}\textstylePolicepardfaut{\textsf{\textcolor{black}{
départ}}}\textstylePolicepardfaut{\textsf{\textbf{\textcolor{black}{.}}}}\textstylePolicepardfaut{\textsf{\textcolor{black}{seconde}}}}

{
\textstylePolicepardfaut{Une variable structurée peut aussi être le
paramètre d’un module, et un module peut également renvoyer une
«~valeur~» de type structuré. Par exemple, l’en-tête d’un module
renvoyant le nombre de secondes écoulées entre une heure de départ et
d’arrivée serait :}}

{\sffamily
\textstyleMotCl{\textcolor{black}{module}}\textstylePolicepardfaut{\textsf{\textcolor{black}{
nbSecondesEcoulées( départ
}}}\textstylePolicepardfaut{{\textsf{\textcolor{black}{↓}}}}\textstylePolicepardfaut{\textsf{\textcolor{black}{,
arrivée
}}}\textstylePolicepardfaut{{\textsf{\textcolor{black}{↓}}}}\textstylePolicepardfaut{\textsf{\textcolor{black}{
: Moment)
}}}\textstylePolicepardfaut{{\textsf{\textcolor{black}{→}}}}\textstylePolicepardfaut{\textsf{\textcolor{black}{
entier}}}}

{
\textstylePolicepardfaut{On pourra aussi lire ou écrire une variable
structurée}\footnote{\textstylePolicepardfaut{Bien que, dans certains
langages, ces opérations devront être décomposées en une lecture ou
écriture de chaque champ de la structure.}}\textstylePolicepardfaut{.}}

{\sffamily
\textstyleMotCl{\textcolor{black}{lire}}\textstylePolicepardfaut{\textsf{\textcolor{black}{
unMoment}}}}

{\sffamily
\textstyleMotCl{\textcolor{black}{écrire}}\textstylePolicepardfaut{\textsf{\textcolor{black}{
unMoment}}}}

{
\textstylePolicepardfaut{Par contre, il n’est pas autorisé d’utiliser
les opérateurs de comparaison ({\textless}, {\textgreater}) pour
comparer des variables structurées (même de même type), car une
relation d’ordre n’accompagne pas toujours les structures utilisées. En
effet, s’il est naturel de vouloir comparer des dates ou des moments,
comment définir une relation d’ordre avec les points du plan ou avec
des cartes d’identités ?}}

{
\textstylePolicepardfaut{Si le besoin de comparer des variables
structurées se fait sentir, il faudra dans ce cas écrire des modules de
comparaison adaptés aux structures utilisées.}}

{
Par facilité, on peut assigner tous les champs en une fois via des
«~\{\}~». Exemple :}

{\sffamily
anniversaire
\textstyleMotCl{{\textcolor{black}{←
}}}\textstyleMotCl{{\textmd{\textcolor{black}{\{1,9,1989\}
\ \ }}}}}

\section{Exemple d’algorithme}
{
Le module ci-dessous reçoit en paramètre deux dates ; la valeur renvoyée
est –1 si la première date est antérieure à la deuxième, 0 si les deux
dates sont identiques ou 1 si la première date est postérieure à la
deuxième.}

{\sffamily
\textstyleMotCl{\textcolor{black}{module}}\textcolor{black}{
comparerDates(date1
}\textstylePolicepardfaut{{\textsf{\textcolor{black}{↓}}}}\textcolor{black}{,
date2
}\textstylePolicepardfaut{{\textsf{\textcolor{black}{↓
:}}}}\textcolor{black}{ Date}\textcolor{black}{)
}\textstylePolicepardfaut{{\textsf{\textcolor{black}{→}}}}\textcolor{black}{
entier}}

{\sffamily
\ \ résultat : entier }

{\sffamily
\textcolor{black}{\ \ résultat
}\textstyleMotCl{{\textcolor{black}{←}}}\textcolor{black}{
–1}}

{\sffamily
\textcolor{black}{\ \ }\textstyleMotCl{\textcolor{black}{si}}\textcolor{black}{
date1}\textbf{\textcolor{black}{.}}\textcolor{black}{année
}\textrm{\textcolor{black}{${\geq}$}}\textcolor{black}{
date2}\textbf{\textcolor{black}{.}}\textcolor{black}{année
}\textstyleMotCl{\textcolor{black}{alors}}}

{\sffamily
\textcolor{black}{\ \ \ \ }\textstyleMotCl{\textcolor{black}{si}}\textcolor{black}{
date1}\textbf{\textcolor{black}{.}}\textcolor{black}{année
{\textgreater}
date2}\textbf{\textcolor{black}{.}}\textcolor{black}{année
}\textstyleMotCl{\textcolor{black}{alors}}}

{\sffamily
\textcolor{black}{\ \ \ \ \ \ résultat
}\textstyleMotCl{{\textcolor{black}{←}}}\textcolor{black}{
1}}

{\sffamily
\ \ \ \ \textstyleMotCl{sinon}\textstyleMotCl{\textmd{ \ // les années
sont identiques}}}

{\sffamily
\textstyleMotCl{\textcolor{black}{\ \ \ \ \ \ si}}\textcolor{black}{
date1}\textbf{\textcolor{black}{.}}\textcolor{black}{mois
}\textrm{\textcolor{black}{${\geq}$}}\textcolor{black}{
date2}\textbf{\textcolor{black}{.}}\textcolor{black}{mois
}\textstyleMotCl{\textcolor{black}{alors}}}

{\sffamily
\textstyleMotCl{\textcolor{black}{\ \ \ \ \ \ \ \ si}}\textcolor{black}{
date1}\textbf{\textcolor{black}{.}}\textcolor{black}{mois
{\textgreater}
date2}\textbf{\textcolor{black}{.}}\textcolor{black}{mois
}\textstyleMotCl{\textcolor{black}{alors}}}

{\sffamily
\textcolor{black}{\ \ \ \ \ \ \ \ \ \ résultat
}\textstyleMotCl{{\textcolor{black}{←}}}\textcolor{black}{
1}}

{\sffamily
\textstyleMotCl{\ \ \ \ \ \ \ \ sinon} // les mois sont identiques}

{\sffamily
\textstyleMotCl{\textcolor{black}{\ \ \ \ \ \ \ \ \ \ si}}\textcolor{black}{
date1}\textbf{\textcolor{black}{.}}\textcolor{black}{jour
}\textrm{\textcolor{black}{${\geq}$}}\textcolor{black}{
date2}\textbf{\textcolor{black}{.}}\textcolor{black}{jour
}\textstyleMotCl{\textcolor{black}{alors}}}

{\sffamily
\textstyleMotCl{\textcolor{black}{\ \ \ \ \ \ \ \ \ \ \ \ si}}\textcolor{black}{
date1}\textbf{\textcolor{black}{.}}\textcolor{black}{jour
{\textgreater}
date2}\textbf{\textcolor{black}{.}}\textcolor{black}{jour
}\textstyleMotCl{\textcolor{black}{alors}}}

{\sffamily
\textcolor{black}{\ \ \ \ \ \ \ \ \ \ \ \ \ \ résultat
}\textstyleMotCl{{\textcolor{black}{←}}}\textcolor{black}{
1}}

{\sffamily
\textstyleMotCl{\ \ \ \ \ \ \ \ \ \ \ \ sinon} // les jours sont
identiques}

{\sffamily
\textcolor{black}{\ \ \ \ \ \ \ \ \ \ \ \ \ \ résultat
}\textstyleMotCl{{\textcolor{black}{←}}}\textcolor{black}{
0}}

{\sffamily
\textstyleMotCl{\ \ \ \ \ \ \ \ \ \ \ \ fin} \textstyleMotCl{si}}

{\sffamily
\textstyleMotCl{\ \ \ \ \ \ \ \ \ \ fin} \textstyleMotCl{si}}

{\sffamily
\textstyleMotCl{\ \ \ \ \ \ \ \ fin} \textstyleMotCl{si}}

{\sffamily
\textstyleMotCl{\ \ \ \ \ \ fin} \textstyleMotCl{si}}

{\sffamily
\textstyleMotCl{\ \ \ \ fin} \textstyleMotCl{si}}

{\sffamily
\textstyleMotCl{\ \ fin} \textstyleMotCl{si}}

{\sffamily
\textstyleMotCl{\ \ retourner} résultat}

{\sffamily
\textstyleMotCl{fin} \textstyleMotCl{module}}

\section{Exercices sur les structures}
\liststyleExercice
\begin{enumerate}
\item {\sffamily\bfseries
Conversion moment-secondes}
\end{enumerate}
{
\textstylePolicepardfaut{Écrire un module qui reçoit en paramètre un
moment d’une journée et qui retourne le nombre de secondes écoulées
depuis minuit jusqu’à ce moment. }}

\liststyleExercice
\setcounter{saveenum}{\value{enumi}}
\begin{enumerate}
\setcounter{enumi}{\value{saveenum}}
\item {\sffamily\bfseries
Conversion secondes-moment}
\end{enumerate}
{
\textstylePolicepardfaut{Écrire un module qui reçoit en paramètre un
nombre de secondes écoulées dans une journée à partir de minuit et qui
retourne le moment correspondant de la journée.}}

\liststyleExercice
\setcounter{saveenum}{\value{enumi}}
\begin{enumerate}
\setcounter{enumi}{\value{saveenum}}
\item {\sffamily\bfseries
Temps écoulé entre 2 moments}
\end{enumerate}
{
\textstylePolicepardfaut{Écrire un module qui reçoit en paramètres deux
moments d’une journée et qui retourne le nombre de secondes séparant
ces deux moments. }}

\liststyleExercice
\setcounter{saveenum}{\value{enumi}}
\begin{enumerate}
\setcounter{enumi}{\value{saveenum}}
\item {\sffamily\bfseries
Milieu de 2 points}
\end{enumerate}
{
\textstylePolicepardfaut{Écrire un module recevant les coordonnées de
deux points distincts du plan et retournant les coordonnées du point
situé au milieu des deux.}}

\liststyleExercice
\setcounter{saveenum}{\value{enumi}}
\begin{enumerate}
\setcounter{enumi}{\value{saveenum}}
\item {\sffamily\bfseries
Distance entre 2 points}
\end{enumerate}
{
\textstylePolicepardfaut{Écrire un module recevant les coordonnées de
deux points distincts du plan et qui }\textstylePolicepardfaut{retourne
la distance entre ces deux points.}}

\liststyleExercice
\setcounter{saveenum}{\value{enumi}}
\begin{enumerate}
\setcounter{enumi}{\value{saveenum}}
\item {\sffamily\bfseries
Un rectangle}
\end{enumerate}
{
\textstylePolicepardfaut{\textcolor{black}{Définir un type
}}\textstylePolicepardfaut{\textbf{\textcolor{black}{Rectangle}}}\textstylePolicepardfaut{\textbf{\textcolor{black}{
}}}\textstylePolicepardfaut{\textcolor{black}{pouvant décrire de façon
commode un rectangle dans un espace à deux dimensions et dont les côtés
sont parallèles aux axes des
}}\textstylePolicepardfaut{\textcolor{black}{coordonnées. Écrire
ensuite}}}

\begin{center}
 [Warning: Image ignored] % Unhandled or unsupported graphics:
%\includegraphics[width=1.134cm,height=1.134cm]{log1-img/log1-img81}

\end{center}
\liststyleNumberingv
\begin{enumerate}
\item {
un module calculant le périmètre d’un rectangle reçu en paramètre ;}
\item {
un module calculant la surface d’un rectangle reçu en paramètre ;}
\item {
un module recevant en paramètre un rectangle R et les coordonnées
d'un point P, et renvoyant la valeur booléenne
\textbf{vrai} si et seulement si le point P est à
l'intérieur du rectangle R ;}
\item {
un module recevant en paramètre un rectangle R et les coordonnées
d'un \ point P, et renvoyant la valeur booléenne
\textbf{vrai} si et seulement si le point P est sur le bord du
rectangle R ;}
\item {
un module recevant en paramètre deux rectangles et renvoyant la valeur
booléenne \textbf{vrai} si et seulement si ces deux rectangles ont une
intersection.}
\end{enumerate}
\liststyleExercice
\begin{enumerate}
\item {\sffamily\bfseries
Un cercle}
\end{enumerate}
{
\textstylePolicepardfaut{\textcolor{black}{Définir un type
}}\textstylePolicepardfaut{\textbf{\textcolor{black}{Cercle}}}\textstylePolicepardfaut{\textbf{\textcolor{black}{
}}}\textstylePolicepardfaut{\textcolor{black}{pouvant décrire de façon
commode un cercle quelconque dans un espace à deux dimensions. Écrire
ensuite}}}

\liststyleNumberingv
\begin{enumerate}
\item {
un module calculant la surface du cercle reçu en paramètre ;}
\item {
un module recevant 2 points en paramètre et retournant le cercle dont le
diamètre est défini par ces 2 points ;}
\item {
un module qui détermine si un point donné est dans un cercle ;}
\item {
un module qui indique si 2 cercles ont une intersection.
}
\end{enumerate}
\liststyleExercice
\begin{enumerate}
\item {\sffamily\bfseries
Validation de date}
\end{enumerate}
{
\textstylePolicepardfaut{\textcolor{black}{Écrire un algorithme qui
valide une date reçue en paramètre sous forme d'une
structure.}}}

